\section{Confluence spectrale, informationnelle et thermodynamique de \(\log3\)}

L’identité

\[
\log3=C_1+C_2-2C_0
=S_{\TRIT}/k_B
=E_P/(k_B\Theta_{\rm clk}),
\]

relie trois réalisations distinctes. Le premier membre est un défaut
résiduel, le deuxième une entropie de décision et le troisième un rapport
entre énergie et température. Leur égalité démontre la conservation du
scalaire sous le changement de réalisation, sans identifier les espaces
d'états.

La première égalité se déduit des trois canaux
résiduels
\[
Z_r(s)=\sum_{\substack{n\ge1\\n\equiv r\ ({\rm mod}\ 3)}}n^{-s},
\qquad r=0,1,2.
\]
Alors
\[
Z_0(s)=3^{-s}\zeta(s),\quad
Z_1(s)+Z_2(s)=(1-3^{-s})\zeta(s),\quad
Z_1(s)-Z_2(s)=L(s,\chi_{-3}).
\]
Si \(C_r\) désigne la partie finie de \(Z_r\) en \(s=1\), le développement
de \(3^{-s}\) donne
\begin{equation}
C_0+C_1+C_2=\gamma,qquad
-2C_0+C_1+C_2=\log3,qquad
C_1-C_2=\frac{\pi}{3\sqrt3}.
\label{eq:residual-three-covectors}
\end{equation}
Les trois covecteurs sont linéairement indépendants. Ils retrouvent, à partir d'une seule
décomposition résiduelle, le terme fini total, la dilatation de la
classe zéro et l'orientation entre les deux classes non nulles. \(\log3\) est
donc un invariant de changement d'échelle avant d'être interprété comme
entropie ou température.

\section{Classification épistémologique des constantes dérivées}

Les contributions du chapitre sont classées selon trois modalités :

\begin{enumerate}
\item attribution d'un opérateur interne à une constante connue, comme dans
le cas de Catalan ;
\item évaluation composée exacte de deux caractères internes, comme
\(L(1,\chi_{12})\) ;
\item obtention d'une constante à partir d'un corps ou d'une famille
construits en interne, comme \(\gamma_{\mathbb Q(\sqrt3)}\) ou \(A_k\).
\end{enumerate}

La provenance mathématique détermine le statut de récupération ou de
nouveauté ; la proximité décimale n'intervient pas dans cette classification.

\begin{remark}[Contrôle négatif]
Le caractère \(\chi_{12}\) doit avoir exactement le motif
\((+,-,-,+)\) dans \((1,5,7,11)\) ; tout autre motif invalide
\cref{eq:L1chi12,eq:L2chi12}. L'identité
\cref{eq:odd-zeta-bendersky} et la factorisation
\cref{eq:dedekind-sqrt3} sont des contrôles indépendants.
\end{remark}

\subsection*{Portée logique et domaine de validité}
\(J^2=-I\), Catalan, \(\chi_{12}\), ses valeurs \(L\), Euler--Kronecker et la tour de Bendersky--Glaisher : \emph{Théorème interne exact}. La réunion en un réseau \(3/4/12\) est une \emph{Formalisation nouvelle} qui conserve les provenances mathématiques.

